\section{Whether a Creature Can See God}

The definition of 1336 asserts something that a considerable weight of Scripture appears to deny outright. Any account of heaven that does not feel that pressure has not read the texts.

\subsection{The texts that say no}

Moses asks to see the divine glory and is refused in terms: ``Thou canst not see my face: for man shall not see me, and live.'' He is placed in a cleft of rock, covered by a hand while the glory passes, and permitted only what is behind: ``thou shalt see my back parts: but my face thou canst not see.''\endnote{Ex 33:20, 23, tracked verse text of the Douay--Rheims (Challoner), read 2026-07-26. The preceding petition is at 33:18, ``Shew me thy glory.''} The Fourth Gospel opens its prologue with a flat negative: ``No man hath seen God at any time: the only begotten Son who is in the Bosom of the Father, he hath declared him.''\endnote{Jn 1:18, tracked verse text, read 2026-07-26.} The First Epistle to Timothy is stronger still, and is not speaking of this life only: God ``only hath immortality and inhabiteth light inaccessible: whom no man hath seen, nor can see.''\endnote{1 Tim 6:16, tracked verse text, read 2026-07-26. \latin{Nec videre potest}: the negative is modal, not merely factual, which is why the verse does most of the work in the objections of both traditions.}

The Greek Fathers pressed these texts as hard as they can be pressed. John Chrysostom, preaching on the prologue, works through every Old Testament theophany---Isaias on the throne, Ezechiel above the cherubim, Daniel's Ancient of Days, Jacob who took his name from seeing God---and disposes of them all: ``all these were instances of (His) condescension, not the vision of the Essence itself unveiled. For had they seen the very Nature, they would not have beheld It under different forms, since that is simple, without form, or parts, or bounding lines.'' Then he extends the denial past the human race: ``what God really is, not only have not the prophets seen, but not even angels nor archangels.'' And he states the principle: ``How can any created nature even see the Uncreated?''\endnote{John Chrysostom, \emph{Homilies on the Gospel of John}, Homily 15 (on Jn 1:18), trans.\ Charles Marriott, \emph{Nicene and Post-Nicene Fathers}, 1st series, vol.~14 (1889), public domain; read 2026-07-26 at New Advent. The New Advent presentation runs scriptural references into the text as inline citations; those are omitted from the quotations here and nothing else is altered. Chrysostom adds that God's essence was invisible even to the angels until the Incarnation: ``even to angels He then became visible, when He put on the Flesh; but before that time they did not so behold Him, because even to them His Essence was invisible.''}

John of Damascus, gathering the Greek tradition three centuries later, makes the same point as a doctrine of God rather than an exegesis: ``what He is in His essence and nature is absolutely incomprehensible and unknowable''; ``God then is infinite and incomprehensible and all that is comprehensible about Him is His infinity and incomprehensibility.''\endnote{John of Damascus, \emph{An Exact Exposition of the Orthodox Faith} I.4, trans.\ S.\ D.\ F.\ Salmond, \emph{Nicene and Post-Nicene Fathers}, 2nd series, vol.~9 (1899), public domain; read 2026-07-26 at New Advent. The chapter heading is ``Concerning the nature of Deity: that it is incomprehensible.''}

It is worth registering how strong this is before going further. These are not marginal voices; the Damascene is the standard dogmatic manual of the Christian East and was quarried by the Latin schoolmen as well. And Aquinas, opening the question in which he will defend the vision of the divine essence, does not soften the objection: he puts Chrysostom first, in his own words---``not only the prophets, but neither the angels nor the archangels have seen that which is God''---and adds Dionysius denying that God is the object of sense, imagination, opinion, reason, or knowledge.\endnote{\emph{Summa theologiae} I, q.~12, a.~1, objections 1 and 3; Latin read 2026-07-26 at Corpus Thomisticum (\latin{Chrysostomus enim, super Ioannem, exponens illud quod dicitur Ioan.\ I, Deum nemo vidit unquam, sic dicit, ipsum quod est Deus, non solum prophetae, sed nec Angeli viderunt nec Archangeli, quod enim creabilis est naturae, qualiter videre poterit quod increabile est?}); English at New Advent.}

\subsection{The texts that say yes}

Against them stands a smaller but very insistent set. ``Blessed are the clean of heart: for they shall see God.''\endnote{Matt 5:8, tracked verse text, read 2026-07-26.} ``We see now through a glass in a dark manner: but then face to face. Now I know in part: but then I shall know even as I am known.''\endnote{1 Cor 13:12, tracked verse text, read 2026-07-26.} ``We know that when he shall appear we shall be like to him: because we shall see him as he is.''\endnote{1 Jn 3:2, tracked verse text, read 2026-07-26. The Douay--Rheims renders the Vulgate's \latin{quoniam} as ``because'' where most modern versions read ``for''; the causal force---likeness because of sight---is what the tradition builds on, and the English Dominican \emph{Summa} translation notes the variant in the text quoted at I, q.~12, a.~5.} And, at the very end of the Apocalypse, in the description of the city: ``And they shall see his face: and his name shall be on their foreheads.''\endnote{Apoc 22:4, tracked verse text, read 2026-07-26.}

The problem is not that some texts are obscure. It is that both sets are clear and they point opposite ways. The whole subsequent history of the doctrine---Latin and Greek---is a history of two different reconciliations, and neither tradition ever simply dropped one side.

\subsection{Aquinas's opening move, and its limits}

The Latin reconciliation is worked out most fully in the twelfth question of the first part of the \emph{Summa theologiae}, thirteen articles on how God is known by creatures. Because this material is usually met in summary, and because summaries of it are unreliable in a specific way---they tend to make the light of glory into a thing seen rather than a power of seeing---it is worth reading in the text.

Article one asks whether any created intellect can see the essence of God. Aquinas begins by conceding the force of the difficulty in his own image: God is supremely knowable in himself, but ``what is supremely knowable in itself, may not be knowable to a particular intellect, on account of the excess of the intelligible object above the intellect; as, for example, the sun, which is supremely visible, cannot be seen by the bat by reason of its excess of light.''\endnote{ST I, q.~12, a.~1, corpus; English at New Advent, Latin at Corpus Thomisticum (\latin{sicut sol, qui est maxime visibilis, videri non potest a vespertilione, propter excessum luminis}), both read 2026-07-26.} Some who considered this, he says, held that no created intellect can ever see God's essence.

He then gives two reasons for rejecting that position, and they are of different kinds, which matters.

The first is theological and he says so: ``as the ultimate beatitude of man consists in the use of his highest function, which is the operation of his intellect; if we suppose that the created intellect could never see God, it would either never attain to beatitude, or its beatitude would consist in something else beside God; \emph{which is opposed to faith}.''\endnote{ST I, q.~12, a.~1, corpus; emphasis added here to mark the appeal. Latin: \latin{Quod est alienum a fide.} Read 2026-07-26 as above.} That is an argument from within revelation, and it presupposes the very thing a sceptic would question.

The second is the famous one: ``there resides in every man a natural desire to know the cause of any effect which he sees; and thence arises wonder in men. But if the intellect of the rational creature could not reach so far as to the first cause of things, the natural desire would remain void.'' Conclusion: ``Hence it must be absolutely granted that the blessed see the essence of God.''\endnote{ST I, q.~12, a.~1, corpus; Latin \latin{Inest enim homini naturale desiderium cognoscendi causam, cum intuetur effectum; et ex hoc admiratio in hominibus consurgit\ldots\ remanebit inane desiderium naturae. Unde simpliciter concedendum est quod beati Dei essentiam videant.} Read 2026-07-26.}

How much that second argument proves has been disputed among Thomists for centuries---whether the desire in question is a properly natural appetite or an elicited one, and whether an argument from it can establish anything about a gratuitous supernatural end without compromising the gratuity. This article names the dispute and does not adjudicate it, because nothing here depends on the outcome: for a Catholic reader the fact of the vision is settled by the definition of 1336, and Aquinas's argument functions as an explanation of a datum rather than as its proof.\endnote{The intra-Thomist controversy over the natural desire for the supernatural is beyond this article's scope and is named, not resolved; see the appendix. What must be preserved either way is the gratuity of the supernatural end, which the First Vatican Council's \emph{Dei Filius} ch.~2 protects by grounding it in God's \latin{infinita bonitas} rather than in any exigency of the creature. That locus is not independently verified for this article and is cited only as a boundary marker.} What the argument does establish, and this is not nothing, is that the Latin tradition thought the vision was owed an explanation---that a creature's arrival at the sight of God is continuous with what a mind is for, rather than an arbitrary prize.

Answering Chrysostom and Dionysius, Aquinas makes the move on which everything afterwards turns: ``Both of these authorities speak of the vision of comprehension.''\endnote{ST I, q.~12, a.~1, reply to objection 1, read 2026-07-26; he adds that Chrysostom himself, immediately after the sentence quoted, says he is speaking ``of the most certain vision of the Father, which is such a perfect consideration and comprehension as the Father has of the Son.'' Whether that reading does justice to Chrysostom's intention is a real question and is taken up in the section on the Eastern answer.} To see is one thing; to comprehend is another. Hold that distinction; the seventh article will cash it.

\subsection{Why no created likeness will do}

Article two asks whether the divine essence is seen through an image or created likeness---a species, in the technical sense, standing in the mind for the thing known, as the likeness of a stone in the eye and not the stone itself. Aquinas gives three reasons why nothing of the sort can serve.

First, on Dionysius's principle, ``by the likeness of a lower order of things, higher things can in no way be known''---no bodily image gives an incorporeal essence, and much less can any created species give the divine one. Second, the essence of God is his very existence, ``which cannot be said of any created form; and so no created form can be the similitude representing the essence of God to the seer.'' Third, the divine essence ``is uncircumscribed, and contains in itself super-eminently whatever can be signified or understood by the created intellect,'' and no created form can represent that, because every created form is determinate---some particular measure of wisdom, or power, or being.\endnote{ST I, q.~12, a.~2, corpus, read 2026-07-26 in English at New Advent and in Latin at Corpus Thomisticum. The scholastic term is \latin{species}, a likeness by which a thing is known; the article's phrase ``created likeness'' renders it.}

Then comes the sentence that decides the shape of the whole doctrine: ``Hence to say that God is seen by some similitude, is to say that the divine essence is not seen at all; which is false.''\endnote{ST I, q.~12, a.~2, corpus; Latin \latin{Unde dicere Deum per similitudinem videri, est dicere divinam essentiam non videri, quod est erroneum.} Read 2026-07-26.} There is no middle option. Either the mind is joined to God himself as the object it knows, or it is not seeing God. Written around 1265--1268, that is precisely the exclusion the constitution of 1336 would promulgate: \latin{nulla mediante creatura in ratione obiecti visi}.

\subsection{Not with the eyes}

Article three settles a question that the devotional imagination almost always gets wrong, and settles it against the imagination. Can the essence of God be seen with the bodily eye? No: ``It is impossible for God to be seen by the sense of sight, or by any other sense, or faculty of the sensitive power. For every such kind of power is the act of a corporeal organ\ldots\ Hence no power of that kind can go beyond corporeal things. For God is incorporeal\ldots\ Hence He cannot be seen by the sense or the imagination, but only by the intellect.''\endnote{ST I, q.~12, a.~3, corpus, read 2026-07-26.}

Job's ``in my flesh I shall see my God'' is therefore not a promise about eyeballs: it means ``that man existing in the flesh after the resurrection will see God.''\endnote{ST I, q.~12, a.~3, reply to objection 1, read 2026-07-26. The Douay--Rheims text of Job 19:26 in the tracked Challoner electronic edition reads ``And I shall be clothed again with my skin, and in my flesh I shall see my God''; the 1899 American edition reads ``I will see my God,'' one of the 1859 verses where the two Douay--Rheims states differ, and the divergence is recorded in the tracked collation table registered with the edition. Neither reading affects the point.} The vision of God is an act of a mind, not of an organ. Every picture in which the blessed \emph{look at} something has already, at that moment, stopped depicting the beatific vision and started depicting something else.

\subsection{The light of glory, read where it is actually written}

Article four asks whether a created intellect can see the divine essence by its own natural powers. The answer turns on a principle Aquinas has used all through the \emph{Summa}: a knower knows according to the mode of its own nature. What exists only in individual matter is connatural to us; separate immaterial natures are connatural to angels; ``to know self-subsistent being is natural to the divine intellect alone.'' Therefore: ``the created intellect cannot see the essence of God, unless God by His grace unites Himself to the created intellect, as an object made intelligible to it.''\endnote{ST I, q.~12, a.~4, corpus, read 2026-07-26. Latin: \latin{Non igitur potest intellectus creatus Deum per essentiam videre, nisi inquantum Deus per suam gratiam se intellectui creato coniungit, ut intelligibile ab ipso.}}

Article five then asks what the creature needs on its own side, and this is where the \latin{lumen gloriae}---the light of glory---appears. The argument is short:

\begin{studyframe}
``Everything which is raised up to what exceeds its nature, must be prepared by some disposition above its nature; as, for example, if air is to receive the form of fire, it must be prepared by some disposition for such a form. But when any created intellect sees the essence of God, the essence of God itself becomes the intelligible form of the intellect. Hence it is necessary that some supernatural disposition should be added to the intellect in order that it may be raised up to such a great and sublime height\ldots\ it is necessary that the power of understanding should be added by divine grace. Now this increase of the intellectual powers is called the illumination of the intellect\ldots\ By this light the blessed are made `deiform'---i.e.\ like to God, according to the saying: `When He shall appear we shall be like to Him, and \ldots\ we shall see Him as He is.'\,''\endnote{ST I, q.~12, a.~5, corpus, read 2026-07-26 in English at New Advent; the ellipsis before ``we shall see Him'' stands where the translators insert the bracketed note that the Vulgate reads ``because'' rather than ``and.'' Latin at Corpus Thomisticum: \latin{oportet quod aliqua dispositio supernaturalis ei superaddatur, ad hoc quod elevetur in tantam sublimitatem\ldots\ Et secundum hoc lumen efficiuntur deiformes, idest Deo similes.} The Scripture cited in the article for the light itself is Apoc 21:23, ``The glory of God hath enlightened it''---glossed by Aquinas as the society of the blessed who see God.}
\end{studyframe}

Two clarifications in the replies do more work than the body of the article, and they are exactly what popular accounts lose.

The first: the created light is not needed to make God intelligible---he is that of himself---but ``in order to enable the intellect to understand in the same way as a habit makes a power abler to act. Even so corporeal light is necessary as regards external sight, inasmuch as it makes the medium actually transparent.''\endnote{ST I, q.~12, a.~5, reply to objection 1, read 2026-07-26.} The second, and decisive: ``This light is required to see the divine essence, not as a similitude in which God is seen, but as a perfection of the intellect, strengthening it to see God. Therefore it may be said that this light is to be described not as a medium in which God is seen, but as one by which He is seen; \emph{and such a medium does not take away the immediate vision of God}.''\endnote{ST I, q.~12, a.~5, reply to objection 2, read 2026-07-26; emphasis added. The scholastic terminology is \latin{medium in quo} (a medium in which, an object seen instead of or alongside) versus \latin{medium quo} (a medium by which, a disposition of the seer). The distinction is the hinge on which the compatibility of the Thomistic account with the 1336 definition turns, and it is also the precise point at which the Eastern tradition declines to follow---for reasons that are not a failure to understand the distinction, as the next section shows.}

That is the whole of it, and it is more modest than the phrase ``light of glory'' usually sounds. The light of glory is not a radiance the blessed look at. It is not a lens they look through. It is a strengthening of the mind itself, a habit-like elevation, so that a created intellect can do what no created intellect can naturally do. The eye is changed; nothing is interposed.

A third reply adds the boundary: this light ``cannot be natural to a creature unless the creature has a divine nature; which is impossible.''\endnote{ST I, q.~12, a.~5, reply to objection 3, read 2026-07-26.} No creature can be born able to see God. Whatever heaven is, it is on this account not the flowering of a capacity we already have.

This is one point at which the Latin school theology touches something more than school theology. The Council of Vienne, in the constitution \emph{Ad nostrum qui} of 1312, listed among the condemned errors of the Beghards and Beguines the proposition ``that every intellectual nature is naturally blessed in itself, and that the soul does not need the light of glory elevating it to see God and to enjoy him blessedly.''\endnote{Council of Vienne, session III, 6 May 1312, constitution \emph{Ad nostrum qui}, fifth listed error, DS 895 (older numbering 475): \latin{Quod quaelibet intellectualis natura in se ipsa naturaliter est beata, quodque anima non indiget lumine gloriae, ipsam elevante ad Deum videndum et eo beate fruendum.} The censure attached to the list (DS 899) is: \latin{Nos sacro approbante Concilio sectam ipsam cum praemissis erroribus damnamus et reprobamus omnino}. Latin read 2026-07-26 at the registered patristica.net presentation; English gloss labelled. The error condemned is a compound proposition, and the condemnation of a listed error of a named sect is not equivalent to a positive definition of the Thomistic account of the \latin{lumen gloriae}; what it does establish is that the denial of any need for an elevating light, coupled with the claim of natural beatitude, is not an open option.} What is condemned there is a compound---natural beatitude plus no need of elevation---and the condemnation of an error is not the canonisation of a particular explanation of the truth it violates. But it does mean that the general shape of the Latin answer is not merely one school's preference.

\subsection{Seeing without comprehending}

Article seven completes the reconciliation, and it is the article that keeps the whole Latin account from arrogance. Do those who see the essence of God comprehend him? Aquinas answers with a sentence of Augustine: ``To comprehend God is impossible for any created intellect; but `for the mind to attain to God in some degree is great beatitude.'\,''\endnote{ST I, q.~12, a.~7, corpus, read 2026-07-26. Latin: \latin{comprehendere Deum impossibile est cuicumque intellectui creato, attingere vero mente Deum qualitercumque, magna est beatitudo, ut dicit Augustinus.} The Augustinian source is the sermon literature; Aquinas cites him without a locus at this point and the article has not independently checked the Augustinian text, which is recorded as a limit.}

The proof is arithmetical in feel. To comprehend a thing is to know it as fully as it is knowable. A thing is knowable in proportion to its actuality; God's being is infinite, so God is infinitely knowable; a created intellect knows the divine essence more or less perfectly according as it receives a greater or lesser light of glory; and a created light received in a created intellect cannot be infinite. ``Hence it is impossible that it should comprehend God.''\endnote{ST I, q.~12, a.~7, corpus, read 2026-07-26; Latin \latin{Cum igitur lumen gloriae creatum, in quocumque intellectu creato receptum, non possit esse infinitum, impossibile est quod aliquis intellectus creatus Deum infinite cognoscat. Unde impossibile est quod Deum comprehendat.}}

Article eight draws the further consequence, which readers of devotional literature will find surprising: the blessed, seeing the divine essence, do \emph{not} see in it everything God does or could do. Effects are seen in a cause to the extent that the cause is comprehended; God is not comprehended; therefore ``no created intellect in seeing God can know all that God does or can do''---though ``of what God does or can do any intellect can know the more, the more perfectly it sees God.''\endnote{ST I, q.~12, a.~8, corpus, read 2026-07-26.} Omniscience is not a side-effect of beatitude. The blessed are not issued the answers.

So the Latin answer, stated whole: the blessed truly see God's own essence, immediately, with no created thing standing in as the object; they do so by a created strengthening of the mind that God gives and no nature possesses; and they never come to the end of him. The infinite remains infinite to them for ever. Whatever else eternity is, on this account it is not a completed inspection.

\subsection{What weight all this carries}

It should be said plainly, because the next sections depend on it: almost none of what has just been laid out is defined doctrine.

Defined: that the blessed see the divine essence, intuitively and facially, with no creature intervening as the object seen (1336); that they see God as he is, unequally, according to merits (1439). Condemned in a compound form: the denial that the soul needs an elevating light (1312). Everything else in this section---the analysis of \latin{species}, the impossibility of a bodily vision, the account of the light of glory as a \latin{medium quo}, the demonstration that comprehension is impossible, the denial that the blessed see all possibles---is the reasoning of a theologian, received with immense authority in the West and with good reason, but reasoning nonetheless.

It is also, in one respect, an unusually honest piece of reasoning, and the honesty is easy to miss. Having asserted that a creature sees God, Aquinas spends the second half of the question saying what that does not mean: not with the eyes, not through an image, not by nature, not comprehensively, not exhaustively, not with all knowledge thrown in. The Latin account of heaven is often accused of intellectualist confidence. Read at the text, it is a chain of denials around a single affirmation---which is roughly what the Christian East says the doctrine of God should look like, and which makes the disagreement between them stranger and more serious than a caricature of either would suggest.
